\chapter[Sermon 73. The Judgment or Punishment of the Purple Whore Is Described: and Also the Sin, and Ungodliness of the Same.]{The Judgment or Punishment of the Purple Whore Is Described: and Also the Sin, and Ungodliness of the Same.}
\label{sermon:073}
\markright{Sermon 73. The Judgment or Punishment of the Purple Whore Is ...}

And there came one of the seven angels, who held the seven vials, and spoke with me, saying to me: Come, I will show you the judgment of the great prostitute, who sits upon many waters, with whom the kings of the earth have committed adultery, and the inhabitants of the earth are drunken with the wine of her fornication. And he carried me away into the wilderness in the spirit. And I saw a woman sitting upon a rose-colored beast, full of names of blasphemy, which had seven heads and ten horns.

\index{Antichrist}\index{Rome}He has hitherto discoursed generally about God’s just judgments; albeit, in the meantime, he has touched on some particular matters concerning Babylon or Rome, rather than reasoned. And now, consequently, individually and evidently, he handles the destruction or end of the beast, and of its image, of old Rome and new—I mean both the empire and papistry—which he seems to point out as it were with the finger. In chapters 13, 14, and 16, he signified something of this; now he pours forth everything with a notable abundance and evidence. And the same thing that I said at the beginning of this passage, I repeat here again: that hereby are affirmed how the justice of God is shown to be the end of good and evil, that the godly are confirmed, and the judgment to come established, as we expressly confess in the Apostles’ Creed. The sum of all is this: old and new Rome, the empire and the Popish kingdom, which is the kingdom of Antichrist, shall perish for sins and great enormities. For he joins together the beast and the image of the beast, the beast and one sitting on the beast, a proud harlot, so that they cannot be separated. Therefore, the passage must be expounded of both empires.

\index{Jerome, Saint}\index{Tertullian}\index{Arethas of Caesarea}And lest anyone should think me to be led by a misguided sentiment in understanding these things concerning Rome, I will demonstrate by the testimony of both God and man that this interpretation is ancient, not new, true, and not contrived. For immediately the angel himself, as we shall hear, explains these things concerning Rome. Arethas, also an ancient interpreter of this book, says that diverse interpreters understand by the whore old Rome. And he continues immediately: While she speaks of the mother of harlotry, whether you will call her old Rome or new, or the time of the coming of Antichrist (behold, he says, the time of Antichrist), you cannot err from the truth, for both cities [Rome and Constantinople] have held the empire, and each has been satiated with the blood of holy martyrs, and so forth. Thus far he. And what the most ancient writer Tertullian and Saint Jerome have spoken repeatedly concerning Babylon and the purple whore, I have reported previously in chapter 14.

And in this order he proceeds. First, he shows the author of this revelation, after he gathers the sum of the revelation or vision. For again he treats by visions, to the end all things might be more lively and evident. Indeed, some divide this work from chapter seven to chapter twenty-one as the sixth vision, as I admonished in the beginning of this work. Then he notes the place and manner of the vision. Finally, he propounds the vision itself, and immediately adjoins the exposition thereof. And in the process of this matter, he uses a judicial kind of pleading, and that after a prophetic manner. For the prophets most often, and in the beginning, set forth the sins and wickedness of the people before the eyes of all men, and then annex unto it the judgment, pain, or punishment. For so does John also at this present.

The author of this horrible vision is the Lord Christ himself, but he uses the ministry of an angel, and that of one who, coming out of the temple of the divine majesty, was appointed with six others to pour out plagues and vials. This is the head minister. It seemed fitting that the judgment of Babylon should be uttered by an angel that had rule over thunderstorms. The Lord Jesus himself will take vengeance on the beast for whom this triumph is reserved. And we understand that such things as are set forth here have proceeded from the high Bishop himself, Jesus Christ, and that the manner of speaking is angelic, heavenly, and godly. Who then shall blame us, if we, using the words of angels and of Christ himself, say that the Bishop of Rome and all Popery is that purple, great, and most common harlot? It appears also to many who seem godly that moderation is neglected when these things are repeated by preachers. It seems that they would shut and stop the mouth of Christ himself. However, they attempt that in vain. For if the preachers hold their peace, the stones will cry. For it behooves that like as the glory of Christ, so the shame of Antichrist should be known to the whole world. But they offend most grievously who, in the sermons made against Antichrist, require I know not what modesty, as though he ought to be spared, which spares no good man. As though that doctrine were not modest, which is taken and received from the mouth and words of Christ. Afterward, in chapter 18, we shall hear the Lord command: “Render unto her as she has rendered unto you,” and so forth.

\index{Papacy}Secondly, he summarizes the essence of everything in a few words and indicates to what we should refer all things. “Come,” says the angel to John, “and I will show you the judgment, condemnation, and punishment of the great whore.” And when he says, “of the great whore,” he intimates what the crime or cause of punishment is: fornication, infidelity, or ungodliness. This vision also pertains to this, so that we might understand how Rome should be punished or destroyed, that is to say, the Roman Empire, or the kingdom of the Pope or of Antichrist, and why or how it deserves to be destroyed. She is a whore, and a great and wandering whore. And who does not know that a marriage is contracted between God and all the faithful? That God is the bridegroom, and the church his spouse? She is then bound and coupled to her husband alone in faith and troth. If she breaks this faith and loves others, gives herself to them, calls upon them, and honors them, she is a whore. Of this I have spoken many times, both in this book and elsewhere.

\index{Ezekiel (Book of)}And it is a very common thing in the Scriptures to call unfaithfulness, impiety, superstition, and idolatry, fornication or whoredom. If anyone desires testimonies of this, he will find them in Judges 8, Isaiah 1, Jeremiah 2 and 3, Ezekiel 16, Hosea 1, 2, and 3, and other places. Rome, therefore, was a great prostitute, and is even now a most stinking harlot, for which she is full of idolatry, the worship of creatures, and abominable superstitions. Neither is she herself only polluted with all filthiness, but compels moreover the whole world to serve, and that to serve in idolatry and superstitions. What will you say that through the wonderful providence of God it came to pass that a woman feigning herself a man did climb up to the See of Rome, was created bishop, and called John 8, which was one Gylberta, a great whore, born at Mentz? For thus would God declare that the Bishop of Rome sits a whore upon the beast. And herein I follow the constant consent of all historians. Nevertheless, I am not ignorant that there are some who have thought how this John was intruded into the seat of an harlot, and for that cause was called an harlot.

Furthermore, old Rome had power to do these things, for she sat upon many waters, that is, had dominion and rule over many people and sundry nations. All the kings of the Earth have committed whoredom with her, lest they have submitted themselves to the Romans, bound themselves in league, and received of them superstitions and idolatry. For the children of Israel were also said to have committed whoredom with the Egyptians, for that they had joined amity with them and were become fellows in profane religions. And so new Rome, the Pope’s kingdom, stretches far and wide, and the kings and princes of the Earth commit whoredom with her. Therefore doth the word of the Lord call it filthy whoredom, which the Romans name a holy bond and obedience. It is added, and they that dwell upon Earth are made drunk. For he signifies that, being infected with errors, yea rather besotted and clean out of their wits, they have been mad in idolatry, and yet rage in their superstitions, like drunkards, and cannot, for fury, receive the preaching of the gospel. Touching this wine of fornication and whoredom, and of that drunkenness, I have spoken in chapter 14. And it is aptly spoken that dwellers upon Earth are made drunk, not so much for that men dwelling upon Earth are made drunken, as for that earthly minds and choked with earthly desires shall become faithful worshippers of the Roman See.

\index{Patmos}Thirdly, he explains the manner of the vision thus: “I was carried away in spirit.” Therefore, while his body remained in Patmos, in spirit he saw a woman sitting on a beast, and consumed with fire. Such are many visions and sights in the prophets. And he notes also the place where he saw the beast, not in heaven, nor in the temple or tabernacle, nor in a fruitful place, but in the wilderness. Isaiah calls the Gentiles and heathens a wilderness. And truly, the Romans and their new superstitions would have had no place in the church, but are outside the church. God forbid that we should acknowledge the church of Rome to be the head of all faithful churches. And at this day, many of those who are called most holy and most reverent differ nothing from the Gentiles, their titles and hypocrisy only excepted, as was spoken of before in chapter 11.

\index{Daniel (Book of)}Fourthly and last, he exhibits this vision, or a glimpse of old and new Rome, and the ruin and destruction of them both, and with all describes most diligently the wickedness of either. First, the beast must be considered, after the woman sitting on the beast. The beast represents the figure of old Rome; the woman, of the new and of Papistry. And the woman sits upon the beast, for the image of the beast has succeeded and has placed her seat in old Rome. For Daniel also affirms that Antichrist shall pitch his seat or palace between two seas, namely, the Adriatic Sea, commonly called the gulf of Venice, and the Tyrrhenian or Tuscan Sea. And the beast is rose-colored, of a red and bright color like crimson, for Rome has been most cruel and bloody, and swimming altogether in the blood of all men, but especially of Christians. How much blood shed Marius, Sylla, Pompey, Julius, and others, as histories Pliny has reported. Rome has with sword and fire destroyed the whole world. The ten persecutions of Christians before the reign of Constantine are most commonly known.

As I showed in the thirteenth chapter, the beast was full of blasphemous names. Rome was filled with chapels and idols, and daily blasphemed God, Christ, and the gospel, tearing the church asunder. Concerning the seven heads and ten horns, there is also speaking in the thirteenth chapter. And certain things will follow in this same chapter, plain enough. And thus much hitherto concerning the old beast: now follows the woman sitting upon the beast.
